\documentclass[12pt]{article}
\usepackage[utf8]{inputenc}
\usepackage[T1]{fontenc}
\usepackage{hyperref}
\title{Huehuetla Tepehua Language Resources for K-12 Curriculum Development}
\author{Compiled by Claude Sonnet 5 (claude-sonnet-5) for Andr\'as Kornai}
\date{}
\begin{document}
\maketitle
\section{Introduction}
Huehuetla Tepehua (Lhiimaqalhqama'; Tepehua de Huehuetla; ISO 639-3 tee) is a Totonacan language of Hidalgo, Mexico. One source reports roughly 1,500--1,600 speakers and a UNESCO status of definitely endangered, lower than the 3,000 in the prompt.
\section{Primary Educational Resources}
Older SIL/ILV primers and storybooks (1948, 1956, 1958) are held in the El Colegio de M\'exico repository; the repository notes their orthography is no longer current: \url{https://repositorio.colmex.mx/concern/books/gt54kp98n?locale=en}, \url{https://repositorio.colmex.mx/concern/books/xd07gv79g?locale=en}, \url{https://repositorio.colmex.mx/concern/books/4q77ft07d?locale=en}. I did not verify which URL corresponds to which title.
\section{Linguistic and Reference Materials}
\emph{Diccionario tepehua: Huehuetla, Hidalgo} (2012, Ediciones Mayahuel; authorship attributed inconsistently): \url{https://openlibrary.org/works/OL43830607W}. Susan Smythe Kung's \emph{A Descriptive Grammar of Huehuetla Tepehua} (2007) is mentioned in search results, without a direct URL found. Omniglot overview: \url{https://omniglot.com/writing/huehuetlatepehua.htm}; numbers: \url{https://www.omniglot.com/language/numbers/huehuetlatepehua.htm}.
\section{Religious and Community Resources}
None found.
\section{Digital and Audio Resources}
Living Dictionary with audio (medicinal plants focus): \url{https://livingdictionaries.app/tepehua-de-huehuetla}. Kung's documentation collection at AILLA: \url{https://islandora-ailla.lib.utexas.edu/islandora/object/ailla%3A124437}.
\section{Curriculum Development Context}
This language is better documented than the others in this batch: a dictionary, a grammar, an archived corpus and an online audio dictionary exist. No school curriculum was found. Community-led work (Lhima'alh'ama' Project) would be the natural partner.
\end{document}